\section{Modelos de lenguaje y adaptación eficiente}
\label{sec:marco_llm}

Un modelo de lenguaje asigna una probabilidad a cada continuación posible de un texto: dado lo escrito
hasta el momento, predice el token siguiente. Con el transformer, estos modelos crecieron en tamaño y en
datos hasta volverse capaces de seguir instrucciones escritas.

\subsection{Grandes modelos de lenguaje}
\label{sec:marco_llm_grandes}

Un \gls{llm} se obtiene por preentrenamiento autosupervisado sobre corpus de miles de millones de
palabras. La tarea es predecir el token siguiente, y la supervisión sale del propio texto, sin etiquetas
manuales. El modelo aprende así la probabilidad de una secuencia como el producto de las probabilidades de
cada token dados los anteriores,
\begin{equation}
    p(x_1, \dots, x_T) = \prod_{t=1}^{T} p(x_t \mid x_1, \dots, x_{t-1}),
    \label{eq:marco_lm}
\end{equation}
y de esa fase sale una distribución de probabilidad sobre el vocabulario condicionada a cualquier
prefijo~\cite{brown2020language,touvron2023llama}.

Según qué parte del transformer conserven, se distinguen dos familias. Los modelos de solo codificador,
como BERT, de 2018, construyen una representación bidireccional del texto, en la que cada token ve a los
anteriores y a los posteriores, y sirven para clasificar y para comparar textos~\cite{devlin2019bert}. Los de
solo decodificador ven únicamente los tokens anteriores y sirven para generar texto. En 2020 se observó que la
pérdida de estos modelos baja de forma predecible al aumentar sus parámetros, sus datos y su
cómputo~\cite{kaplan2020scaling}, y GPT-3, con 175 mil millones de parámetros, mostró las capacidades que
aparecen con esa escala~\cite{brown2020language}. Desde 2023 hay además familias de pesos abiertos, que se
pueden descargar y ajustar, como LLaMA~\cite{touvron2023llama} y Mistral, cuyo modelo de 7 mil millones de
parámetros supera a uno de LLaMA~2 de 13 mil millones~\cite{jiang2023mistral}.

La familia Phi siguió el camino opuesto al aumento de escala: reducir el tamaño y mejorar los datos. Sus
primeros modelos se entrenaron con textos de calidad de libro de texto, en parte
sintéticos~\cite{gunasekar2023textbooks}, y Phi-3-mini, con 3800 millones de parámetros, rinde a la par de
modelos mucho mayores y cabe en un teléfono. Su variante Phi-3.5-vision, de 4200 millones de parámetros,
agrega entrada visual~\cite{abdin2024phi3}. Todos estos son modelos de solo decodificador.

El comportamiento del modelo se ajusta después del preentrenamiento. El ajuste fino supervisado lo entrena
sobre pares instrucción--respuesta, y la alineación con \gls{rlhf} lo entrena a preferir las
respuestas que un evaluador humano marca como mejores. Un modelo alineado responde a instrucciones, de modo
que se le puede pedir una tarea escribiéndola en lenguaje natural~\cite{ouyang2022training}.

Todo el texto que el modelo recibe antes de generar, la instrucción y cualquier información adicional,
forma su contexto y se llama \emph{prompt}. Como la probabilidad de la Ecuación~\ref{eq:marco_lm} está
condicionada en ese prefijo, cambiar el \emph{prompt} cambia lo que el modelo genera sin modificar sus
pesos. Los modelos grandes aprovechan además la información y los ejemplos incluidos en el \emph{prompt}, sin
ningún paso de entrenamiento, lo que se conoce como aprendizaje en contexto
(\emph{in-context learning})~\cite{brown2020language}.

Dado un \emph{prompt}, el texto se genera token a token: en cada paso el modelo devuelve una distribución
sobre el vocabulario, se elige un token y se lo agrega al prefijo. La regla con que se elige es la
estrategia de decodificación. La decodificación voraz (\emph{greedy decoding}) toma siempre el token más
probable, y en la generación abierta, como los otros métodos que buscan el texto más probable, tiende a
repetirse~\cite{holtzman2020curious}. Los
métodos de muestreo (\emph{sampling}) sortean el token siguiente según la distribución del modelo. La
temperatura $T$ divide los logits antes de la softmax,
\begin{equation}
    p_i = \frac{\exp(z_i / T)}{\sum_{j} \exp(z_j / T)},
    \label{eq:marco_temperatura}
\end{equation}
de modo que con $T<1$ la probabilidad se concentra en los tokens más probables y con $T>1$ se reparte; es
la misma operación que la Sección~\ref{sec:marco_eval_calib} usa para calibrar un clasificador. El muestreo de
núcleo (\emph{top-p}) sortea entre el conjunto más chico de tokens cuya probabilidad acumulada alcanza un
umbral $p$, de modo que la cantidad de candidatos crece cuando el modelo duda y se achica cuando está seguro.
Con eso evita tanto la repetición de la decodificación voraz como la incoherencia de muestrear de la
distribución completa~\cite{holtzman2020curious}. Como la estrategia de decodificación cambia el texto
generado, forma parte de la configuración de un generador y hay que declararla al comparar modelos.

\subsection{Ajuste eficiente de parámetros: LoRA}
\label{sec:marco_llm_lora}

Adaptar uno de estos modelos a un dominio con el ajuste fino completo de la
Sección~\ref{sec:marco_transferencia} tropieza con la memoria: con Adam, cada parámetro entrenado arrastra su
gradiente y dos estados adicionales, y para miles de millones de parámetros eso puede superar la memoria de una
GPU. Los métodos de \gls{peft} congelan el modelo preentrenado y entrenan solo un conjunto pequeño de
parámetros, y reducen así la memoria que el ajuste necesita~\cite{hu2021lora}.

\Gls{lora}, propuesta en 2021, es uno de ellos. Parte de que el cambio que el ajuste fino introduce en una
matriz de pesos tiene rango bajo, es decir, se puede escribir como el producto de dos matrices delgadas. Sobre
cada matriz congelada $W$ que se elige adaptar entrena entonces un incremento de bajo rango,
\begin{equation}
    W' = W + \frac{\alpha}{r}\,BA, \qquad B \in \mathbb{R}^{d\times r},\quad A \in \mathbb{R}^{r\times k},\quad
    r \ll \min(d, k),
    \label{eq:marco_lora}
\end{equation}
y solo actualiza $A$ y $B$ (Figura~\ref{fig:marco_lora}). $A$ se inicializa con valores aleatorios y $B$ en
cero, de modo que el incremento vale cero al empezar y el modelo parte exactamente del preentrenado. El rango
$r$ fija cuántos parámetros tiene el incremento y con ello su capacidad de adaptación, y el factor $\alpha$
escala su efecto, de modo que se puede cambiar $r$ sin volver a ajustar la tasa de
aprendizaje~\cite{hu2021lora}.

% Figura: reparametrización de LoRA. Adaptada de hu2021lora (Figura 1).
% Se usa en sec:marco_llm_lora, después de la Ecuación de LoRA.

\begin{figure}[H]
    \centering
    \begin{tikzpicture}[x=1mm, y=1mm]

        \node[figentrada, text width=10mm, minimum height=7mm] (x) at (0, 0) {$\mathbf{x}$};

        % Rama congelada.
        \node[fignodo, text width=36mm, minimum height=24mm, fill=figrelleno] (W) at (-24, 25)
            {pesos preentrenados\\$W \in \mathbb{R}^{d\times k}$\\[2pt]{\footnotesize congelados}};

        % Rama entrenable: A reduce a r, B vuelve a d.
        \draw[draw=fignaranja!70!black, fill=fignaranja!25, thick] (8, 16) -- (32, 16) -- (26, 22) -- (14, 22) -- cycle;
        \node[font=\footnotesize] at (20, 19) {$A \in \mathbb{R}^{r\times k}$};
        \draw[draw=fignaranja!70!black, fill=fignaranja!25, thick] (14, 28) -- (26, 28) -- (32, 34) -- (8, 34) -- cycle;
        \node[font=\footnotesize] at (20, 31) {$B \in \mathbb{R}^{d\times r}$};
        \node[font=\footnotesize, text=fignaranja!70!black, align=left, anchor=west] at (34, 25)
            {entrenables\\rango $r \ll \min(d,k)$};

        % Suma y salida.
        \node[circle, draw=figtrazo, inner sep=1pt, font=\small] (sum) at (0, 46) {$+$};
        \node[figsalida, text width=26mm, minimum height=7mm] (h) at (0, 58) {$\mathbf{h} = W\mathbf{x} + BA\mathbf{x}$};

        \draw[figflecha] (x.north) -- ++(0, 4) -| (W.south);
        \draw[figflecha] (x.north) -- ++(0, 4) -| (20, 16);
        \draw[figflecha] (20, 22) -- (20, 28);
        \draw[figflecha] (W.north) |- (sum.west);
        \draw[figflecha] (20, 34) |- (sum.east);
        \draw[figflecha] (sum.north) -- (h.south);

    \end{tikzpicture}
    \caption{Reparametrización de LoRA con el incremento de bajo rango $BA$. Adaptada de \textcite{hu2021lora}.}
    \label{fig:marco_lora}
\end{figure}


El método no fija qué matrices adaptar. El trabajo original lo aplicó solo a las de la atención
(Sección~\ref{sec:marco_atencion}), en la mayoría de los experimentos a las proyecciones $Q$ y $V$, y dejó las
del perceptrón de cada bloque como trabajo futuro, aunque señala que se puede aplicar a cualquier subconjunto de
matrices. Cuantas más matrices se adaptan y mayor es
$r$, más parámetros se entrenan. Aun así son una fracción mínima: en GPT-3 redujo los parámetros entrenables
unas 10\,000 veces y la memoria de la GPU unas tres veces~\cite{hu2021lora}. Al terminar, el incremento se suma
a $W$ y el modelo resultante se ejecuta igual que el original, sin latencia adicional.

Con un \gls{llm} adaptable en una GPU de consumo, falta conectarlo con la imagen. Eso exige llevar los rasgos
de la imagen al espacio de tokens del modelo de lenguaje, o bien alinear imagen y texto en un espacio común.
